\section{Discussion}
\label{sec:discussion}

% Subsections still to be written by their owners:
%   - Main findings on SAE features, probes, and interventions
%   - Alternative explanations for the probe and intervention results

\subsection{Limits of the behavioral labels}
\label{sec:disc-labels}

The probe targets for the combined collection are generated by a language
model. On the Adoption--Retention decision they agree with the frozen
rule-based labels in every sampled case, and on the team-reviewed episodes
they match the team's judgment more often than the rules do. They remain
model judgments, however, so we report them as LLM-derived labels rather
than ground truth. The two sources differ mainly on ambiguous replies that
the rules place in Out. A probe result that holds under one label source but
not the other should be read with that difference in mind, and the frozen
rule-based labels remain available for checking any headline result on the
original collection.

Both label sources depend on the Critic's recommended answer, the only
feedback the Solver sees in the candidate presentation. Neither judges
whether an answer is factually correct, so the binary target measures
uptake of feedback, not improvement in accuracy.

\subsection{Limits of the synthetic pilot}
\label{sec:disc-synthetic}

The pilot has 100 questions per group, which is enough to detect the large
difference in disagreement but leaves wide intervals for any single hop type.
Its synthetic questions use fictional entities, so the models cannot draw on
prior knowledge in either direction. This makes them a cleaner test of
reading the passages, but less like real questions. The comparison with
MuSiQue is also affected by MuSiQue's own defects and by strict answer
matching, both of which inflate MuSiQue's apparent difficulty. We therefore
treat the pilot as evidence about which generation choices make questions
too easy, not as a measure of how much harder synthetic questions should be.
